\section{Good models, linear triviality, and boundary consequences}
\label{sec:consequences}

The real-boundary induction and the final line-bundle descent have
different natural scopes. We first record good-model existence over
\(\C\), including the precise smooth canonical comparison used by
fiber-space applications. We then deduce actual rational-linear
triviality in Iitaka dimension zero over every field covered by
Theorem~\ref{thm:main}.
Finally, over \(\C\), normalization gluing and a section-extension
theorem give two consequences for reducible boundaries and nonnormal pairs.

\begin{theorem}[Real-boundary good models over \(\C\)]
\label{thm:good-models-complex}
Every projective log canonical pair \((X,B)\) over \(\C\), with
effective real boundary and pseudo-effective real Cartier adjoint
\(K_X+B\), has a good log minimal model. Real semiampleness has the
meaning fixed in Section~\ref{sec:induction}.
\end{theorem}
\begin{proof}
The simultaneous induction in Section~\ref{sec:completion} starts with
dimension zero and proves the full real-boundary conclusion of
Proposition~\ref{prop:inductive-reduction} at each dimension. In
particular, its conclusion applies to the given pair before any
rational-boundary or nefness specialization.
\end{proof}

\begin{corollary}[Smooth canonical good models]\label{cor:smooth-good-model}
Let \(X\) be a smooth projective complex variety with pseudo-effective
\(K_X\). There is a normal projective \(\Q\)-factorial klt variety
\(V\) and a \(K_X\)-negative birational contraction \(X\dashrightarrow V\)
such that \(K_V\) is \(\Q\)-Cartier and semiample. On a common smooth
resolution \(p:W\to X\), \(q:W\to V\), compatible canonical divisors
satisfy
\[
 p^*K_X=q^*K_V+E,
 \qquad E\geq0\quad\text{and }E\text{ is }q\text{-exceptional}.
\]
In particular \(\kappa(X,K_X)=\kappa(V,K_V)\).
\end{corollary}
\begin{proof}
Apply Theorem~\ref{thm:good-models-complex} to the klt pair \((X,0)\).
Choose its \(\Q\)-factorial log minimal model, as in the construction
of Proposition~\ref{prop:inductive-reduction}. The klt clause of
Lemma~\ref{mmp:comparison} shows that it extracts no divisor, remains
klt, and has the displayed effective exceptional comparison.
Its boundary is zero, being the strict transform of the zero boundary
with no extracted divisors.
Discrepancies strictly improve at contracted divisors, which is the
\(K_X\)-negative condition; semiampleness is the good-model conclusion.
For sufficiently divisible \(m>0\), the effective \(q\)-exceptional
divisor \(mE\) satisfies \(q_*\OO_W(mE)=\OO_V\), since \(V\) is
normal. Projection formula and the displayed comparison identify the
pluricanonical section spaces. Their ratios agree in the common
function field, proving the assertion about Iitaka dimension.
\end{proof}

\begin{lemma}\label{lem:semiample-zero}
Let \(Z\) be a nonempty integral projective variety over an algebraically
closed field, and let \(A\) be a semiample \(\Q\)-Cartier divisor.
If \(\kappa(Z,A)=0\), then \(A\sim_{\Q}0\).
\end{lemma}
\begin{proof}
Choose a positive integer \(m\) such that the line bundle
\(L=\OO_Z(mA)\) is globally generated. Its complete system gives
\[
 \phi:Z\longrightarrow\PP(H^0(Z,L)^*),\qquad
 L\simeq\phi^*\OO(1).
\]
The image has dimension at most \(\kappa(Z,A)=0\), so it is a point.
A linear form nonzero at that point pulls back to a nowhere-vanishing
section of \(L\). Thus \(L\simeq\OO_Z\), an actual line-bundle
isomorphism, and \(A\sim_{\Q}0\).
\end{proof}

\begin{corollary}\label{cor:qlinear-zero}
Let \((X,B)\) satisfy the hypotheses of Theorem~\ref{thm:main}.
If \(\kappa(X,K_X+B)=0\), then \(K_X+B\sim_{\Q}0\).
In particular, every nef rational adjoint of Iitaka dimension zero
on a projective klt fivefold over an algebraically closed field of
characteristic zero is \(\Q\)-linearly trivial.
\end{corollary}
\begin{proof}
Theorem~\ref{thm:main} gives semiampleness on the original normal
variety. Lemma~\ref{lem:semiample-zero} gives a positive Cartier
multiple with trivial associated line bundle. Klt pairs are lc, so
the fivefold assertion is a special case.
\end{proof}

Thus the nef adjoint also has numerical dimension zero. This is a
consequence of the trivializing Cartier multiple, not an additional
hypothesis. The multiple may depend on the pair; no uniform index is
asserted.

\subsection{Semi-log-canonical abundance and dlt extension}

We use the standard semi-log-canonical convention, in which the underlying
projective scheme may be reducible or nonnormal. It is reduced, satisfies
\(S_2\), is pure-dimensional, and has ordinary double crossings in
codimension one. No component of the effective boundary is contained in the
codimension-one singular locus. If \(\nu:X^\nu\to X\) is the normalization,
the boundary \(\Theta\)
defined by
\[
 K_{X^\nu}+\Theta=\nu^*(K_X+\Delta)
\]
includes the conductor boundary, and \((X^\nu,\Theta)\) is componentwise
log canonical.
These are the conventions of \cite[Definition~2.5]{FujinoGongyo2014}.

\begin{corollary}[Semi-log-canonical abundance over \(\C\)]
\label{cor:slc-abundance}
Let \((X,\Delta)\) be a projective semi-log-canonical pair over \(\C\),
where \(\Delta\) is an effective rational divisor and \(K_X+\Delta\) is
\(\Q\)-Cartier. If \(K_X+\Delta\) is nef, then it is semiample.
\end{corollary}
\begin{proof}
Write \(X^\nu=\coprod_i X_i^\nu\) for the normalization and put
\(\Theta_i=\Theta|_{X_i^\nu}\). Each \((X_i^\nu,\Theta_i)\) is a normal
projective log canonical pair with effective rational boundary. The adjoint
\(K_{X_i^\nu}+\Theta_i\) is the pullback of \(K_X+\Delta\), so it is nef.
Theorem~\ref{thm:main} makes each of these adjoints semiample. There are only
finitely many components, hence one common positive Cartier multiple is
globally generated on their disjoint union. Thus \(\nu^*(K_X+\Delta)\) is
semiample. The normalization gluing theorem
\cite[Theorem~1.5]{FujinoGongyo2014} now gives semiampleness on \(X\).
\end{proof}

\begin{corollary}[Supported extension for dlt pairs over \(\C\)]
\label{cor:dlt-extension}
Let \((X,\Delta)\) be a projective dlt pair over \(\C\), where \(\Delta\)
is an effective rational divisor, and put \(S=\lfloor\Delta\rfloor\).
Assume that \(K_X+\Delta\) is nef and that, for an effective rational
divisor \(D\),
\[
 K_X+\Delta\sim_{\Q}D\geq0,
 \qquad S\subseteq\Supp D.
\]
Then, for every integer \(m\geq2\) for which \(m(K_X+\Delta)\) is Cartier,
the restriction map
\[
 H^0\bigl(X,\OO_X(m(K_X+\Delta))\bigr)
 \longrightarrow
 H^0\bigl(S,\OO_X(m(K_X+\Delta))|_S\bigr)
\]
is surjective.
\end{corollary}
\begin{proof}
Set \(B=\Delta-S\). The pair \((X,S+B)\) is dlt, \(B\) is an effective
rational divisor, and \(\lfloor S+B\rfloor=S\). Since a dlt pair is log
canonical, Theorem~\ref{thm:main} makes its nef adjoint semiample.
The pair and the displayed effective representative therefore satisfy the
hypotheses of \cite[Proposition~5.12]{FujinoGongyo2014}, which gives the
stated restriction surjectivity for every Cartier multiple with \(m\geq2\).
This is the supported extension formulation in
\cite[Conjecture~5.8]{FujinoGongyo2014}.
\end{proof}

The first corollary settles the rational-boundary form of the semi-log-canonical
abundance conjecture \cite[Conjecture~1.3]{FujinoGongyo2014} in every
dimension over \(\C\). The second settles the supported dlt extension
conjecture \cite[Conjectures~1.10 and~5.8]{FujinoGongyo2014} with the precise
Cartier-multiple range supplied by Proposition~5.12. Both are applications
of the cited implications after normal abundance has been established here.
The support condition \(S\subseteq\Supp D\) is retained; no arbitrary
restriction-surjectivity statement is asserted.
